\originalchapter{Higgs 机制与质量的检验}\label{ch:higgs}
\originalsection{标量真空与规范场质量}
在规范拉格朗日量中直接加入 $m_A^2A_\mu A^\mu/2$ 通常破坏规范结构. Higgs 机制引入与规范场耦合的标量, 通过非零真空值生成质量. 对最简单 U(1) 模型,
$\mathcal L=|D_\mu\phi|^2-\mu^2|\phi|^2-\lambda|\phi|^4-\tfrac14F^2$, $\lambda>0$, $\mu^2<0$. 极小化得 $|\phi|^2=v^2/2=-\mu^2/(2\lambda)$. 写 $\phi=(v+h)\e^{\ii\chi/v}/\sqrt2$, 协变导数产生
\begin{equation}
 |D_\mu\phi|^2=\frac12(\partial h)^2+\frac12(v+h)^2(\partial_\mu\chi/v-gA_\mu)^2.
\end{equation}
单位规范把相位吸收到规范场纵向自由度, 得 $M_A=gv$, $m_h^2=2\lambda v^2$. 规范冗余本身没有成为被实验打破的新守恒定律; ``自发破缺''是选定真空和表示下的常用语言.

\originalsection{电弱质量关系}
标准模型 Higgs 双重态取 $Y=1$, 真空值 $(0,v/\sqrt2)^T$. 展开协变导数, 带电组合 $W^\pm=(W^1\mp\ii W^2)/\sqrt2$ 质量为 $gv/2$, 中性质量矩阵有一个零特征值与一个非零特征值:
\begin{equation}
 M_W=\frac{gv}2,\qquad M_Z=\frac{v\sqrt{g^2+g'^2}}2,
 \qquad M_W=M_Z\cos\theta_W.
\end{equation}
零特征值对应光子, 所以剩余 U(1)$_{\rm EM}$ 保留无质量光子. 树级关系有辐射修正, 不能把精密实验差异直接说成理论全部失败.

\originalsection{Yukawa 耦合}
以电子为例, $-y_e\bar L\phi e_R+\mathrm{h.c.}$ 在真空后给 $m_e=y_ev/\sqrt2$ 与 $h\bar ee$ 耦合 $m_e/v$. 夸克和带电轻子的质量由各自 Yukawa 系数给出; 理论没有用这一个机制独自预测所有 Yukawa 数值. 质子质量主要来自 QCD 的束缚与场动力学, 不是把三个轻夸克 Higgs 质量相加.

对 Higgs 衰变到 Dirac 费米子,
\begin{equation}
 \Gamma(h\to f\bar f)=N_c\frac{m_f^2m_h}{8\pi v^2}
 \left(1-\frac{4m_f^2}{m_h^2}\right)^{3/2}.
\end{equation}
颜色因子对夸克为 $3$, 对轻子为 $1$. 题目若只保留 $b,c,\tau,\mu$ 四道, 分支比可按 $N_cm_f^2$ 归一化; 它不是实际 Higgs 的完整分支比, 实际还有 $WW^*,ZZ^*,gg,\gamma\gamma$ 等通道, 夸克质量还须指定尺度.

\originalsection{产生、衰变与发现}
强子对撞机可经 gluon fusion、矢量玻色子融合、伴随 W/Z 产生或与 top 对伴生来产生 Higgs. Gluon fusion 和 $h\to\gamma\gamma$ 虽然末态无直接树级 Yukawa 耦合, 可经带电或带色粒子回路发生. 在狭宽近似中产生率与分支比相乘, 观测还需探测效率.

重建 $\gamma\gamma$ 或四轻子不变质量可寻找局部共振, 同时检验背景、质量、角分布与多个通道一致性. 课程材料以 2020 年的实验图景说明这些方法; 本稿不复制受限实验协作组照片和图, 不把当年的未来测量规划写成现在已经完成的结果. Higgs 质量、耦合和自相互作用各是不同测量, ``发现一个新粒子''不等于所有理论参数都已确定.
